%% 56장 모수 검정
\chapter{모수 검정}
\label{ch:parametric-tests}

\abstract{정규 모형을 바탕으로 하는 고전적인 검정들을 다룬다. 평균에 관한 $t$ 검정은 52장의 결과 덕분에 귀무 분포가 정확하지만 두 표본의 분산이 다르면 합동 분산을 쓰는 Student의 검정이 크게 틀린다. 크기가 작은 쪽의 표준편차가 세 배(분산이 아홉 배)이면 명목 수준 $5\%$의 검정이 $26\%$로 기각하는 것을 모의실험으로 확인하고 Welch--Satterthwaite 근사 자유도를 쓰는 Welch 검정이 이 왜곡을 바로잡는 까닭을 유도한다. 일원 분산분석의 제곱합 분해를 증명하고 분산이 같은지를 보는 검정들을 견준다. 정규성 아래의 우도비 검정을 고쳐 쓴 Bartlett 검정은 $t_3$ 분포에서 $41\%$, 지수 분포에서 $38\%$로 잘못 기각한다. 4차 적률이 유한하면 통계량의 극한이 카이제곱이 아니라 카이제곱을 초과 첨도 $\gamma_2$에 따라 $1+\gamma_2/2$배로 늘인 분포여서 생기는 일임을 본다. 끝으로 Shapiro--Wilk와 적률에 기초한 정규성 검정을 정리한다.}

\section{평균의 $t$ 검정}
\label{sec:t-tests}

\subsection{한 표본과 짝지은 표본}

$X_1,\dots,X_n$이 독립이고 $N(\mu,\sigma^2)$을 따르면 52장에서 $T=(\bar X-\mu)/(S/\sqrt n)\sim t_{n-1}$임을 보였다. 그래서 $H_0:\mu=\mu_0$의 검정은 $T_0=(\bar X-\mu_0)/(S/\sqrt n)$을 $t_{n-1}$의 분위수와 견주고 55장의 쌍대성으로 구간 $\bar X\pm t_{n-1,1-\alpha/2}S/\sqrt n$을 얻는다. 같은 개체를 두 번 잰 짝지은 자료는 차이 $D_i=X_i-Y_i$에 한 표본 검정을 쓴다. 짝 안의 두 측정이 양의 상관을 이루면 개체 사이의 변동이 차이에서 빠지므로 두 표본을 따로 다룰 때보다 검정력이 크다. 상관이 $0$에 가까우면 자유도만 $2n-2$에서 $n-1$로 줄어 검정력이 오히려 조금 작다.

\subsection{두 표본: Student와 Welch}

독립인 두 표본 $X_1,\dots,X_{n_1}\sim N(\mu_1,\sigma_1^2)$, $Y_1,\dots,Y_{n_2}\sim N(\mu_2,\sigma_2^2)$에서 $H_0:\mu_1=\mu_2$를 검정하자. 분산이 같다고 가정하면($\sigma_1=\sigma_2$) 합동 분산
\[
S_p^2=\frac{(n_1-1)S_1^2+(n_2-1)S_2^2}{n_1+n_2-2}
\]
으로 만든 $T=(\bar X-\bar Y)/(S_p\sqrt{1/n_1+1/n_2})$가 정확히 $t_{n_1+n_2-2}$를 따른다. $(n_1+n_2-2)S_p^2/\sigma^2$이 독립인 카이제곱의 합이어서 $\chi^2_{n_1+n_2-2}$이고 표본평균들과 독립이기 때문이다. 이 통계량으로 하는 검정이 Student의 두 표본 검정이다.

분산이 다르면 사정이 다르다. 평균 차의 분산은 $\sigma_1^2/n_1+\sigma_2^2/n_2$인데 $S_p^2(1/n_1+1/n_2)$는 각 표본의 분산을 자유도 $n_i-1$로 가중하므로 크기가 큰 쪽의 분산에 끌린다. 크기가 작은 쪽의 분산이 크면 평균 차의 분산을 과소평가해 검정이 지나치게 자주 기각하고 반대면 지나치게 드물게 기각한다. 모의실험(반복 $4$만 번)으로 확인하면 \cref{tab:welch-student}과 같다. $n_1=10$, $n_2=40$이고 작은 쪽의 표준편차가 $3$배이면 Student 검정은 명목 수준 $5\%$에서 $25.8\%$로 기각하고 큰 쪽이 $3$배이면 $0.17\%$로만 기각한다. 표본 크기가 같으면 두 검정이 비슷하다.

\begin{table}[t]
\begin{sdtable}{colspec={ccccc}, row{1}={font=\small\bfseries}}
\toprule
$n_1$ & $n_2$ & $\sigma_1/\sigma_2$ & Student & Welch \\
\midrule
$10$ & $10$ & $1$ & $0.051$ & $0.050$ \\
$10$ & $40$ & $1$ & $0.049$ & $0.051$ \\
$10$ & $40$ & $3$ & $0.258$ & $0.050$ \\
$10$ & $40$ & $1/3$ & $0.002$ & $0.048$ \\
$25$ & $25$ & $1/3$ & $0.054$ & $0.051$ \\
\bottomrule
\end{sdtable}
\caption{귀무가설 아래 두 표본 $t$ 검정의 실제 기각률(명목 수준 $0.05$, 반복 $4$만 번). 두 정규 모집단의 평균은 같고 표준편차의 비만 다르다.}
\label{tab:welch-student}
\end{table}

Welch는 각 표본의 분산을 따로 쓴 통계량
\[
T_W=\frac{\bar X-\bar Y}{\sqrt{S_1^2/n_1+S_2^2/n_2}}
\]
을 제안했다\cite{welch1938,welch1947}. 분모의 제곱 $V=S_1^2/n_1+S_2^2/n_2$은 서로 다른 척도의 카이제곱의 합이어서 정확한 분포가 간단하지 않으므로 적률을 맞추어 $V$를 $c\chi^2_\nu/\nu$ 꼴로 근사한다. 이 근사는 Welch의 1938년 논문에 처음 나왔고 Satterthwaite가 분산 성분 일반으로 넓혔으므로\cite{satterthwaite1946} 그 자유도를 Welch--Satterthwaite 자유도라 부른다. $a_i=\sigma_i^2/n_i$로 쓰면 $E[V]=a_1+a_2$이고 $\operatorname{Var}V=2a_1^2/(n_1-1)+2a_2^2/(n_2-1)$이다($\operatorname{Var}S_i^2=2\sigma_i^4/(n_i-1)$). $c\chi^2_\nu/\nu$의 평균 $c$와 분산 $2c^2/\nu$를 이 두 값에 맞추면 $c=a_1+a_2$이고
\[
\nu=\frac{(a_1+a_2)^2}{a_1^2/(n_1-1)+a_2^2/(n_2-1)}
\]
이다. 실제로는 $a_i$에 $S_i^2/n_i$를 넣는다. 이 자유도는 일반적으로 정수가 아니며 $\min(n_1,n_2)-1$과 $n_1+n_2-2$ 사이에 있다. \cref{tab:welch-student}에서 Welch 검정은 모든 경우에 명목 수준을 지킨다. 분산이 같은 경우에도 잃는 것이 거의 없으므로 분산이 같다는 확신이 없으면 Welch 검정을 쓰는 것이 낫다.

\section{일원 분산분석}
\label{sec:anova}

$k$개 집단에서 각각 $n_j$개의 관측 $Y_{ij}$가 독립이고 $N(\mu_j,\sigma^2)$을 따를 때 $H_0:\mu_1=\dots=\mu_k$를 검정하는 것이 \term{일원 분산분석}{one-way analysis of variance}이다. 전체 평균을 $\bar Y$, 집단 평균을 $\bar Y_j$, 전체 관측 수를 $N$이라 하자.

\begin{proposition}
\label{prop:anova-decomposition}
$\sum_{j,i}(Y_{ij}-\bar Y)^2=\sum_jn_j(\bar Y_j-\bar Y)^2+\sum_{j,i}(Y_{ij}-\bar Y_j)^2$이다.
\end{proposition}

\begin{proof}
$Y_{ij}-\bar Y=(\bar Y_j-\bar Y)+(Y_{ij}-\bar Y_j)$를 제곱해 더하면 교차항이 $2\sum_j(\bar Y_j-\bar Y)\sum_i(Y_{ij}-\bar Y_j)$이고 안쪽 합은 집단 평균의 정의로 $0$이다.
\end{proof}

세 항을 전체, 집단 사이(SSB), 집단 안(SSW)의 제곱합이라 한다. 52장의 직교 사영 논법으로 귀무가설 아래 $\mathrm{SSB}/\sigma^2\sim\chi^2_{k-1}$, $\mathrm{SSW}/\sigma^2\sim\chi^2_{N-k}$이고 둘은 독립이므로
\[
F=\frac{\mathrm{SSB}/(k-1)}{\mathrm{SSW}/(N-k)}\sim F(k-1,N-k)
\]
이다. 평균이 다르면 SSB가 커지므로 $F$가 클 때 기각한다. $k=2$이면 $F$는 Student 두 표본 통계량의 제곱이다. 이 검정도 Student 검정처럼 등분산을 가정하므로 분산이 다르고 집단 크기가 다르면 같은 문제가 생긴다. Welch의 분산분석은 집단마다 무게 $w_j=n_j/S_j^2$를 주어 같은 생각을 $k$개 집단으로 넓힌다\cite{welch1951}.

\section{분산이 같은지의 검정}
\label{sec:variance-homogeneity}

분산분석의 가정을 확인하거나 분산 자체를 비교하려면 $H_0:\sigma_1^2=\dots=\sigma_k^2$을 검정한다. 정규 모형의 우도비 검정을 자유도로 고쳐 쓴 것이 \term{Bartlett 검정}{Bartlett's test}이다\cite{bartlett1937}. 통계량은 합동 분산과 집단 분산의 로그의 차
\[
B=\frac{(N-k)\log S_p^2-\sum_j(n_j-1)\log S_j^2}{1+\frac1{3(k-1)}\bigl(\sum_j\frac1{n_j-1}-\frac1{N-k}\bigr)}
\]
이고 귀무가설 아래 근사적으로 $\chi^2_{k-1}$이다. 분모는 근사를 개선하는 보정 계수이다.

Bartlett 검정은 정규성에 극도로 민감하다. $S_j^2$의 분산이 정규 모집단에서는 $2\sigma^4/(n_j-1)$이지만 일반적으로는 $\sigma^4(2+\gamma_2)/n_j$쯤이다($\gamma_2$는 초과 첨도). 그래서 $B$의 극한은 $\chi^2_{k-1}$이 아니라 $(1+\gamma_2/2)\chi^2_{k-1}$이다. 이 극한은 4차 적률이 유한해야 성립한다. 4차 적률이 없는 $t_3$에서는 $B$가 표본 크기와 함께 커져 기각률이 $1$로 간다. 지수 분포($\gamma_2=6$)이면 $4\chi^2_{k-1}$이어서 $k=3$에서 점근적으로 $P(4\chi^2_2>5.99)=e^{-0.749}=0.47$의 확률로 잘못 기각한다.

이 민감성을 피하려는 검정들이 있다. \term{Levene 검정}{Levene's test}은 편차의 절댓값 $Z_{ij}=|Y_{ij}-\tilde Y_j|$($\tilde Y_j$는 $j$번째 집단의 중심이고 Levene은 평균을 썼다)에 일원 분산분석의 $F$ 검정을 쓴다\cite{levene1960}. 분산이 큰 집단일수록 편차의 절댓값도 평균이 커서 이 평균을 견주는 것이 곧 분산을 견주는 셈이 된다. Brown과 Forsythe는 중심 $\tilde Y_j$를 평균 대신 중앙값으로 두면 치우친 분포에서 더 강건하다는 것을 보였다\cite{brown1974}. scipy도 중앙값을 기본 중심으로 쓴다. \term{Fligner--Killeen 검정}{Fligner--Killeen test}은 집단 중앙값에서 잰 절댓값 편차에 순위를 매기고 그 순위를 반정규 분포의 분위수로 바꾸어 쓴다. 두 표본 검정\cite{fligner1976}을 Conover, Johnson, Johnson이 $k$개 집단으로 넓힌 꼴이다\cite{conover1981}.

$k=3$, 집단마다 $n=20$에서 세 검정의 귀무가설 아래 기각률을 모의실험(반복 $8000$번)으로 견주면 \cref{tab:variance-tests}와 같다. 정규 분포에서는 셋 다 명목 수준 근처이지만 자유도 $3$인 $t$ 분포에서 Bartlett 검정은 $41\%$, 지수 분포에서 $38\%$로 기각한다. 중앙값 중심의 Levene 검정은 세 분포 모두에서 명목 수준 아래에 머문다. Fligner--Killeen 검정은 지수 분포에서 $8.8\%$로 명목 수준을 넘는다.

\begin{table}[t]
\begin{sdtable}{colspec={lccc}, row{1}={font=\small\bfseries}}
\toprule
분포 & Bartlett & Levene(중앙값) & Fligner--Killeen \\
\midrule
정규 & $0.051$ & $0.038$ & $0.034$ \\
$t_3$ & $0.413$ & $0.038$ & $0.041$ \\
지수 & $0.381$ & $0.047$ & $0.088$ \\
\bottomrule
\end{sdtable}
\caption{세 집단($n=20$씩)의 분산이 같을 때 등분산 검정의 실제 기각률(명목 수준 $0.05$, 반복 $8000$번).}
\label{tab:variance-tests}
\end{table}

\section{정규성 검정}
\label{sec:normality}

모수 검정의 정규성 가정을 확인하는 검정도 여럿이다. \term{Shapiro--Wilk 검정}{Shapiro--Wilk test}의 통계량은
\[
W=\frac{\bigl(\sum_ia_iX_{(i)}\bigr)^2}{\sum_i(X_i-\bar X)^2}
\]
이다\cite{shapiro1965}. 무게는 표준 정규 표본의 순서통계량의 기댓값 벡터 $m$과 공분산 행렬 $V$로 $a=V^{-1}m/\|V^{-1}m\|$이다. $a$는 합이 $0$이고 길이가 $1$이므로 $W$는 순서통계량 $X_{(i)}$와 무게 $a_i$의 표본 상관의 제곱이다. 정규 표본에서는 순서통계량이 정규 점수와 거의 일직선이어서 $W$가 $1$에 가깝다. 무게와 $W$의 귀무 분포는 닫힌 꼴이 없어 Royston의 근사 알고리즘(AS R94)으로 계산하며\cite{royston1995} scipy와 R도 이 알고리즘을 쓴다. 적률에 기초한 검정도 있다. 표본 왜도와 첨도를 각각 정규 근사로 표준화해 제곱합을 만든 D'Agostino--Pearson의 $K^2$\cite{dagostino1973}과, 더 단순한
\[
\mathrm{JB}=n\Bigl(\frac{g_1^2}6+\frac{g_2^2}{24}\Bigr)
\]
(Jarque--Bera\cite{jarque1980})이 대표적이다. 둘 다 근사적으로 $\chi^2_2$이지만 Jarque--Bera 쪽은 근사의 수렴이 느리다. 명목 수준 $5\%$에서는 작은 표본에서 보수적이지만(정규 표본 $n=20$에서 실제 기각률이 약 $2.7\%$) $1\%$ 수준에서는 표본이 수백이어도 명목보다 자주 기각한다. 정규성 검정은 표본이 작으면 검정력이 낮고 크면 실제로 무시해도 될 작은 벗어남까지 기각한다. 그러니 그 결과를 모수 검정을 쓸지 정하는 기계적인 기준으로 삼기보다 Q--Q 그림과 함께 보는 것이 좋다.

\begin{instardust}[\code{docs/design/htest.md} §1.1 H19--H38]
\code{sd.htest}의 $t$ 검정은 scipy의 이름(\code{ttest\_1samp}, \code{ttest\_ind}, \code{ttest\_rel})을 따르고 \code{ttest\_ind(equal\_var=False)}가 Welch 검정이다. 분산분석은 \code{f\_oneway}이고 \code{equal\_var=False}로 Welch의 분산분석을 한다. 자유도는 Welch처럼 정수가 아닐 수 있으므로 어디서나 float이고 $F$ 꼴 검정은 \code{df\_num}과 \code{df\_denom}을 결과에 둔다. 표준오차나 분산이 $0$인 퇴화한 입력의 통계량과 $p$값은 무한대가 아니라 NaN이다. 분산이 $0$인지는 기술통계의 \code{skew}와 같은 판정($m_2\le(\varepsilon\cdot\text{평균})^2$) 하나로 정한다. Shapiro--Wilk는 Royston의 AS R94를 따른다. Jarque--Bera의 커널은 htest 한곳에 두었으므로 회귀의 잔차 진단 \code{r.jarque\_bera()}도 같은 비트의 결과를 낸다.
\end{instardust}

\section{맺음}

정규 모형의 고전 검정들은 52장의 표본 분포 결과로 정확한 귀무 분포를 얻는다. 그러나 정확성은 가정에 기대고 가정이 어긋나면 결과가 크게 틀릴 수 있다. 두 표본의 분산이 다를 때 Student 검정이, 꼬리가 무거울 때 Bartlett 검정이 그렇다. Welch의 근사 자유도와 Levene류의 강건한 검정은 가정을 약하게 해 이 위험을 줄이며 잃는 검정력은 대개 작다. 다음 장에서는 분포의 모양을 전혀 가정하지 않는 순위 검정과 경험분포에 기초한 검정으로 넘어간다.

\chapterreferences
